
\documentclass{article}
\usepackage{colortbl}
\usepackage{makecell}
\usepackage{multirow}
\usepackage{supertabular}

\begin{document}

\newcounter{utterance}

\centering \large Interaction Transcript for game `codenames', experiment `association\_difficult', episode 3 with reflexion{-}llama{-}3.1{-}8b{-}instruct.
\vspace{24pt}

{ \footnotesize  \setcounter{utterance}{1}
\setlength{\tabcolsep}{0pt}
\begin{supertabular}{c@{$\;$}|p{.15\linewidth}@{}p{.15\linewidth}p{.15\linewidth}p{.15\linewidth}p{.15\linewidth}p{.15\linewidth}}
    \# & $\;$A & \multicolumn{4}{c}{Game Master} & $\;\:$B\\
    \hline

    \theutterance \stepcounter{utterance}  
    & & & \multicolumn{2}{p{0.3\linewidth}}{
        \cellcolor[rgb]{0.9,0.9,0.9}{
            \makecell[{{p{\linewidth}}}]{
                \texttt{\tiny{[GM$|$GM]}}
                {'hidden': {'team': ['radish', 'apple', 'anatomy'], 'innocent': ['parsley', 'foundation', 'green'], 'opponent': [], 'assassin': ['door']}, 'revealed': {'team': {'team': [], 'innocent': [], 'opponent': [], 'assassin': []}, 'opponent': {'team': [], 'innocent': [], 'opponent': [], 'assassin': []}}}
            }
        }
    }
    & & \\ \\

    \theutterance \stepcounter{utterance}  
    & & \multicolumn{4}{p{0.6\linewidth}}{
        \cellcolor[rgb]{0.9,0.9,0.9}{
            \makecell[{{p{\linewidth}}}]{
                \texttt{\tiny{[P1$\langle$GM]}}
                \texttt{You have attempted this same task before and did not succeed. These are the plans you wrote after those attempts.} \\
\texttt{Trial \#1: New plan: } \\
\\ 
\texttt{In my previous attempts, I focused on providing clues that were related to multiple words, but I didn't consider the potential for my teammate to guess an opponent's word or a distractor word. This led to my teammate's turn ending prematurely, which ultimately resulted in my team losing the game.} \\
\\ 
\texttt{Next time, I will aim to provide clues that are more specific and targeted towards my team's words, while also being mindful of the potential for my teammate to guess an opponent's word or a distractor word. I will also try to provide clues that are more challenging for my teammate to guess, but not so challenging that they are unable to make a correct guess.} \\
\\ 
\texttt{To achieve this, I will focus on providing clues that are related to a single word or a small group of words, and I will try to use more abstract and creative connections between the clue word and the target words. I will also pay closer attention to the words that have been removed from the lists and try to avoid providing clues that are too similar to those words.} \\
\\ 
\texttt{For example, instead of providing a clue like "CLUE: plant TARGETS: radish, apple, garlic, mesa", I might provide a clue like "CLUE: sprout TARGETS: radish, garlic" or "CLUE: fruit TARGETS: apple, radish". This will allow my teammate to make more targeted guesses and increase the chances of my team winning the game.} \\
\\ 
\texttt{Trial \#2: New plan: } \\
\\ 
\texttt{In my previous attempts, I focused on providing clues that were related to multiple words, but I didn't consider the potential for my teammate to guess an opponent's word or a distractor word. This led to my teammate's turn ending prematurely, which ultimately resulted in my team losing the game.} \\
\\ 
\texttt{Next time, I will aim to provide clues that are more specific and targeted towards my team's words, while also being mindful of the potential for my teammate to guess an opponent's word or a distractor word. I will also try to provide clues that are more challenging for my teammate to guess, but not so challenging that they are unable to make a correct guess.} \\
\\ 
\texttt{To achieve this, I will focus on providing clues that are related to a single word or a small group of words, and I will try to use more abstract and creative connections between the clue word and the target words. I will also pay closer attention to the words that have been removed from the lists and try to avoid providing clues that are too similar to those words.} \\
\\ 
\texttt{For example, instead of providing a clue like "CLUE: plant TARGETS: radish, apple, garlic, mesa", I might provide a clue like "CLUE: sprout TARGETS: radish, garlic" or "CLUE: fruit TARGETS: apple, radish". This will allow my teammate to make more targeted guesses and increase the chances of my team winning the game.} \\
\\ 
\texttt{In this specific case, I will focus on providing a clue that is more closely related to the word "anatomy". Since "anatomy" is the only word that is not a type of food, I will try to provide a clue that is more abstract and related to the concept of "anatomy". For example, I might provide a clue like "CLUE: study TARGETS: anatomy" or "CLUE: human TARGETS: anatomy". This will allow my teammate to make a more targeted guess and increase the chances of my team winning the game.} \\
\\ 
\texttt{Trial \#3: New plan: } \\
\\ 
\texttt{In my previous attempts, I focused on providing clues that were related to multiple words, but I didn't consider the potential for my teammate to guess an opponent's word or a distractor word. This led to my teammate's turn ending prematurely, which ultimately resulted in my team losing the game.} \\
\\ 
\texttt{Next time, I will aim to provide clues that are more specific and targeted towards my team's words, while also being mindful of the potential for my teammate to guess an opponent's word or a distractor word. I will also try to provide clues that are more challenging for my teammate to guess, but not so challenging that they are unable to make a correct guess.} \\
\\ 
\texttt{To achieve this, I will focus on providing clues that are related to a single word or a small group of words, and I will try to use more abstract and creative connections between the clue word and the target words. I will also pay closer attention to the words that have been removed from the lists and try to avoid providing clues that are too similar to those words.} \\
\\ 
\texttt{In this specific case, I will focus on providing a clue that is more closely related to the word "anatomy". Since "anatomy" is the only word that is not a type of food, I will try to provide a clue that is more abstract and related to the concept of "anatomy". For example, I might provide a clue like "CLUE: study TARGETS: anatomy" or "CLUE: human TARGETS: anatomy". This will allow my teammate to make a more targeted guess and increase the chances of my team winning the game.} \\
\\ 
\texttt{However, I also need to consider the words "radish" and "apple", which are both types of food. To provide a clue that is more closely related to these words, I might use a word that is related to the concept of food or eating. For example, I might provide a clue like "CLUE: eat TARGETS: radish, apple" or "CLUE: salad TARGETS: radish, apple". This will allow my teammate to make a more targeted guess and increase the chances of my team winning the game.} \\
\\ 
\texttt{Ultimately, my goal is to provide clues that are challenging but not impossible for my teammate to guess, and that will allow my team to win the game.} \\
\\ 
\\ 
\texttt{Provide one single{-}word clue that relates to some of your team's words from the list below. You can choose to how many and to which words your clue relates to.} \\
\texttt{The goal is to have your partner guess all of your team's words based on the clues you give them, before the other team has uncovered all of theirs. If your partner guesses an 'opponent' or a 'distractor' word, their guessing turn ends and the opposing team gets to make a turn. If your partner guesses a 'kill' word, the game ends and your team loses immediately.} \\
\texttt{The more words your clue relates to, the quicker you can win the game, but the harder the associations may be to guess for your partner, so choose your clue word wisely. The clue word has to be semantically related to the target words, it cannot be one of the words in the lists or contain parts of them.} \\
\\ 
\texttt{Always give your single{-}word clue and your comma{-}separated list of related target words in the following format and make your answers as short as possible, never include any other text than is required in this form:} \\
\\ 
\texttt{CLUE: <WORD>} \\
\texttt{TARGETS: <TARGETS>} \\
\\ 
\texttt{Your team words are: radish, apple, anatomy.} \\
\texttt{Your opponent's team words are: .} \\
\texttt{Distractor words are: parsley, foundation, green.} \\
\texttt{The kill words are: door.} \\
            }
        }
    }
    & \\ \\

    \theutterance \stepcounter{utterance}  
    & \multicolumn{4}{p{0.6\linewidth}}{
        \cellcolor[rgb]{0.9,0.9,0.9}{
            \makecell[{{p{\linewidth}}}]{
                \texttt{\tiny{[P1$\rangle$GM]}}
                \texttt{CLUE: salad} \\
\texttt{TARGETS: radish, apple} \\
            }
        }
    }
    & & \\ \\

    \theutterance \stepcounter{utterance}  
    & & & \multicolumn{2}{p{0.3\linewidth}}{
        \cellcolor[rgb]{0.9,0.9,0.9}{
            \makecell[{{p{\linewidth}}}]{
                \texttt{\tiny{[GM$|$GM]}}
                \texttt{salad} \\
            }
        }
    }
    & & \\ \\

    \theutterance \stepcounter{utterance}  
    & & & \multicolumn{2}{p{0.3\linewidth}}{
        \cellcolor[rgb]{0.9,0.9,0.9}{
            \makecell[{{p{\linewidth}}}]{
                \texttt{\tiny{[GM$|$GM]}}
                ['radish', 'apple']
            }
        }
    }
    & & \\ \\

    \theutterance \stepcounter{utterance}  
    & & & \multicolumn{2}{p{0.3\linewidth}}{
        \cellcolor[rgb]{0.9,0.9,0.9}{
            \makecell[{{p{\linewidth}}}]{
                \texttt{\tiny{[GM$|$GM]}}
                {'word': 'radish', 'assignment': 'team'}
            }
        }
    }
    & & \\ \\

    \theutterance \stepcounter{utterance}  
    & & & \multicolumn{2}{p{0.3\linewidth}}{
        \cellcolor[rgb]{0.9,0.9,0.9}{
            \makecell[{{p{\linewidth}}}]{
                \texttt{\tiny{[GM$|$GM]}}
                {'word': 'apple', 'assignment': 'team'}
            }
        }
    }
    & & \\ \\

    \theutterance \stepcounter{utterance}  
    & & & \multicolumn{2}{p{0.3\linewidth}}{
        \cellcolor[rgb]{0.9,0.9,0.9}{
            \makecell[{{p{\linewidth}}}]{
                \texttt{\tiny{[GM$|$GM]}}
                \texttt{opponent has won} \\
            }
        }
    }
    & & \\ \\

    \theutterance \stepcounter{utterance}  
    & \multicolumn{4}{p{0.6\linewidth}}{
        \cellcolor[rgb]{0.95,0.95,0.95}{
            \makecell[{{p{\linewidth}}}]{
                \texttt{\tiny{[P2$\rightarrow$P1]}}
                \texttt{Post{-}episode Reflexion (not sent to GM):} \\
\texttt{New plan: } \\
\\ 
\texttt{In my previous attempts, I focused on providing clues that were related to multiple words, but I didn't consider the potential for my teammate to guess an opponent's word or a distractor word. This led to my teammate's turn ending prematurely, which ultimately resulted in my team losing the game.} \\
\\ 
\texttt{Next time, I will aim to provide clues that are more specific and targeted towards my team's words, while also being mindful of the potential for my teammate to guess an opponent's word or a distractor word. I will also try to provide clues that are more challenging for my teammate to guess, but not so challenging that they are unable to make a correct guess.} \\
\\ 
\texttt{To achieve this, I will focus on providing clues that are related to a single word or a small group of words, and I will try to use more abstract and creative connections between the clue word and the target words. I will also pay closer attention to the words that have been removed from the lists and try to avoid providing clues that are too similar to those words.} \\
\\ 
\texttt{In this specific case, I will focus on providing a clue that is more closely related to the word "anatomy". Since "anatomy" is the only word that is not a type of food, I will try to provide a clue that is more abstract and related to the concept of "anatomy". For example, I might provide a clue like "CLUE: study TARGETS: anatomy" or "CLUE: human TARGETS: anatomy". This will allow my teammate to make a more targeted guess and increase the chances of my team winning the game.} \\
\\ 
\texttt{However, I also need to consider the words "radish" and "apple", which are both types of food. To provide a clue that is more closely related to these words, I might use a word that is related to the concept of food or eating. For example, I might provide a clue like "CLUE: eat TARGETS: radish, apple" or "CLUE: salad TARGETS: radish, apple". This will allow my teammate to make a more targeted guess and increase the chances of my team winning the game.} \\
\\ 
\texttt{Ultimately, my goal is to provide clues that are challenging but not impossible for my teammate to guess, and that will allow my team to win the game.} \\
\\ 
\texttt{CLUE: study} \\
\texttt{TARGETS: anatomy} \\
            }
        }
    }
    & & \\ \\

\end{supertabular}
}

\end{document}
